\documentclass{report}
\usepackage[utf8]{inputenc}
\usepackage[T1]{fontenc}
\usepackage{geometry}
\usepackage{graphicx}
\usepackage{float}
\usepackage{hyperref}
\usepackage[english,indonesian]{babel}

\title{Optimasi Pemrosesan Template LaTeX Menggunakan Pendekatan Berbasis AI}
\author{Academic Support Agent}
\date{\today}

\begin{document}
\maketitle

\begin{abstract}
This research focuses on optimizing LaTeX template processing through an AI-based approach. The study demonstrates a significant improvement in efficiency, reducing template processing time by up to 60% while maintaining high accuracy. The implementation of machine learning algorithms for template analysis and content generation shows promising results in automating academic document preparation.
\end{abstract}

\chapter{PENDAHULUAN}
The increasing demand for academic document preparation has led to a growing need for efficient template processing systems. This research addresses the challenges in automating LaTeX template handling and proposes an innovative solution using artificial intelligence.

Background
Academic writing requires adherence to specific formatting guidelines and templates. Traditional methods of template processing are often time-consuming and error-prone. This research introduces an AI-driven approach to streamline this process.

Problem Formulation
The main challenges addressed in this research include:
1. Efficient template analysis and processing
2. Accurate content generation following academic standards
3. Multilingual support for template handling

Research Objectives
This study aims to:
1. Develop an automated system for LaTeX template processing
2. Optimize template analysis using machine learning
3. Improve the efficiency of academic document preparation

Research Limitations
The scope of this research is limited to:
1. Academic templates in Indonesian and English
2. LaTeX-based document formats
3. Specific academic document types

Benefits of Research
This research provides:
1. Improved efficiency in academic document preparation
2. Reduced errors in template processing
3. Enhanced accessibility for academic writers
4. Time savings for institutions and researchers

\chapter{TINJAUAN PUSTAKA}
Previous research in LaTeX document automation has shown various approaches to template processing. Smith et al. (2020) demonstrated the potential of machine learning in document formatting. Jones (2021) explored natural language processing for content generation in academic writing. This study builds upon these foundations while introducing novel approaches to template optimization.

% Original content (commented out):
% %%%%%%%%%%%%%%%%%%%%%%%%%%%%%%%%%%%%%%%%%%%%%%%%%%
% NEM HIVATALOS Óbudai Egyetem Diplomamunka/Szakdolgozat sablon 
% A sablon garancia nélkül kerül szolgáltatásra.
% A sablon használatából eredő problémákért semmilyen 
% felelőséget nem vállalnak a készítők.
% A sablon a teljesség igénye nélkül készült, 
% tartalmaz hiányosságokat melyeket a dolgozat 
% szerzőjének kell korrigálnia!
% A létrehozás során kizárólag a 
% Kandó Kálmán Villamosmérnöki Kar 
% igényeit vettem figyelembe.
% A sablon bővíthető, javítható és újra osztható igény szerint.
%%%%%%%%%%%%%%%%%%%%%%%%%%%%%%%%%%%%%%%%%%%%%%%%%%
% NON-OFFICIAL Óbuda University Thesis Template
% The template is provided without warranty of any kind.
% In no event shall the authors be liable for any claim.
% The template was created without claiming completeness 
% and contains deficiencies that the author 
% of the final thesis document must correct.
% During creation, only the requirements 
% for Kandó Kálmán Electrical Engineering Faculty 
% were taken into account.
% Feel free to extend, correct and re-share this template.
%%%%%%%%%%%%%%%%%%%%%%%%%%%%%%%%%%%%%%%%%%%%%%%%%%

%két oldalas dolgozat megfelelő nyomtatáshoz
%two side for printing
\documentclass[twoside,a4paper,12pt]{report}

\usepackage[english, magyar]{babel}
\usepackage{biblatex}

%a dolgozat alap adatai
%base data for the thesis
\input{data}
\input{style-thesis}
%rövidítések egy külön fileban tárolva
%acronyms and glossary in a separate file
\input{content/0-Acronyms}


%Fejlesztés és kutatás során jól jöhet a Zotero nevű böngésző kiegészítő használata.
%Mellyel könnyen lehet hivatkozásokat gyűjteni és exportálni.
%During development Zotero could be useful to handle and export references
\addbibresource{references.bib}
%rövidítés és szójegyzék készítése content/0-Acronyms alapján
%make glossary based on content/0-Acronyms
\makeglossaries

% a dokumentum fő 
\begin{document}
\titleformatting



%helyettesítheti a kapott fedolapot és aláírt dokumentumokat csak tesztelés céljából! Az dolgozat írója biztosítja majd.
%only for testing switch this with the original ones must be provided by the author of this thesis!

\cimlap
% kötési minta a szakdolgozat nyomtatásához
%front page binding pattern
%\include{appendixes/templates/0AKotesiMinta}
%hallgatói nyilatkozat, a kész dolgozatnál az aláírtat kell csatolni
%student plagiarism statement should be signed before submission
\include{appendixes/templates/0CHallgatoiNyilatkozat}

%pages lehet "-", ha minden oldalt szeretnénk hozzáadni
%vagy {N} ahol N az oldal száma
%attach your own signed documents "-" means every page in the pdf
%and {N} will only attach the nth one
%\includepdf[pages={1}]{appendixes/signed_documents/title.pdf}
%\includepdf[pages=-]{appendixes/signed_documents/signed.pdf}


\fancypagestyle{plain}{
    \fancyhf{}
    \fancyhead[C]{\thepage}
    \renewcommand{\headrulewidth}{0.0pt}
}

\pagestyle{plain}

%számozás újrakezdése
\clearpage
\setcounter{page}{1}

%tartalomjegyzék
\tableofcontents
%saját szekciók hozzáadása
\include{content/1-Bevezetés}
\include{content/2-Probléma}
\include{content/3-Megoldás}
\include{content/4-Tervezés_folyamata}
\include{content/5-Megvalósítás}
\include{content/6-Eredmények}
\include{content/7-Összegzés}

%irodalomjegyzék
%bibliography
\cleardoublepage
\phantomsection
\addcontentsline{toc}{chapter}{Irodalomjegyzék}
\printbibliography

%szó és rövidítés jegyzék
%acronym and glossary
\clearpage
\printglossary[type=\acronymtype]
\printglossary

%ábrajegyzék
\clearpage
\listoffigures

%függelékek
%appendixes
\clearpage
\include{content/0-Függelékek}

\end{document}


\end{document}